\chapter{Diseño de los Datos de Entrada}

\section{Propósito del capítulo}

En la solución propuesta, los datos de entrada constituyen el vínculo operativo entre el espacio físico del museo, el dispositivo del visitante y el conocimiento curatorial que sustenta la mediación digital. Por ello, este capítulo no se limita a enumerar tipos de datos, sino que define los insumos mínimos que la solución necesita para reconocer contexto de visita, interpretar consultas y producir respuestas coherentes con la sala o zona en la que se encuentra el usuario.

La relevancia de este capítulo dentro de la tesis es metodológica. Antes de describir la arquitectura y el despliegue de la solución, resulta necesario establecer qué señales, contenidos, configuraciones y metadatos alimentan su funcionamiento. Esta especificación permite que el diseño posterior sea trazable a entradas concretas y verificables, en lugar de apoyarse en supuestos implícitos. En consecuencia, el capítulo organiza los datos de entrada de la solución propuesta según su procedencia, función y nivel de estructuración, y fija criterios mínimos de calidad para su uso posterior en la solución propuesta.

Además, en el estado actual del proyecto ya existen tres repositorios que materializan esta lógica de entrada en un prototipo distribuido: \texttt{iot-museiq}, centrado en la emisión BLE desde ESP32; \texttt{iot}, que implementa la aplicación móvil y su base local de recorrido; y \texttt{museRAG}, que expone la capa conversacional y la base vectorial. Por ello, el presente capítulo no solo define un contrato de datos deseable en términos conceptuales, sino que también explicita cómo dicho contrato ya empieza a concretarse en artefactos de software desarrollados para la tesis.

\section{Clasificación general de los datos de entrada}

Desde una perspectiva de ingeniería aplicada, los datos de entrada de la solución propuesta pueden agruparse en cuatro categorías principales. La primera corresponde a los datos de contexto físico-espacial, compuestos por identificadores BLE y observaciones de intensidad de señal recibida en el smartphone. La segunda agrupa datos de orientación y movimiento derivados de sensores del dispositivo móvil. La tercera comprende las entradas directas del visitante mediante voz, texto o eventos de interfaz. La cuarta está constituida por el corpus curatorial y sus metadatos, que alimentan la capa conversacional basada en recuperación aumentada.

Esta clasificación es coherente con la base tecnológica del proyecto presentada en el capítulo anterior. BLE y los sensores móviles permiten reconocer contexto situado de manera funcional; la interacción multimodal habilita consultas accesibles durante el recorrido; y la capa de inteligencia requiere un corpus institucional estructurado para sostener respuestas ancladas en evidencia documental~\cite{leitch2023indoor,android2026positionsensors,android2026speech,android2026tts,lewis2020rag}. A ello se suma una distinción temporal importante: mientras las señales sensoriales y eventos de interacción se generan en tiempo real, el corpus curatorial y la configuración espacial corresponden a datos relativamente estables que deben mantenerse bajo control institucional. En el prototipo ya implementado, esta separación también es visible a nivel computacional: \texttt{iot-museiq} produce datos efímeros de advertising BLE, \texttt{iot} mantiene datos curatoriales y de recorrido en almacenamiento local, y \texttt{museRAG} consolida fuentes documentales persistentes para recuperación semántica. La Figura~\ref{fig:clasificacion-datos-entrada-museiq} resume esta organización general.

\begin{figure}[!ht]
    \centering
    \includegraphics[width=0.92\textwidth]{clasificacion_datos_entrada_museiq.png}
    \caption[Clasificación de los datos de entrada de la solución propuesta]{Clasificación general de los datos de entrada requeridos por la solución propuesta. Fuente: elaboración propia.}
    \label{fig:clasificacion-datos-entrada-museiq}
\end{figure}

La clasificación anterior puede comprenderse mejor si se observa el flujo completo que siguen los datos dentro del prototipo. En la práctica, la solución propuesta integra señales físicas del entorno, contexto de recorrido y conocimiento curatorial en una secuencia continua que culmina en una respuesta contextualizada para el visitante, tal como se aprecia en la Figura~\ref{fig:flujo-datos-museiq-ilustrativo}.

\begin{figure}[!ht]
    \centering
    \includegraphics[width=0.98\textwidth]{flujo-datos-museiq-ilustrativo.png}
    \caption[Flujo general de datos en la solución propuesta]{Representación ilustrativa del flujo general de datos en la solución propuesta, desde la captura de señales físicas y consultas del visitante hasta la recuperación curatorial y la entrega de la respuesta contextual. Fuente: elaboración propia con apoyo de generación visual.}
    \label{fig:flujo-datos-museiq-ilustrativo}
\end{figure}

\section{Datos de entrada para reconocimiento espacial por sala o zona}

\subsection{Identificadores BLE y estructura de beacon}

La localización funcional propuesta para la solución propuesta se apoya en una infraestructura de beacons BLE basada en ESP32. Cada beacon transmite un identificador detectable por el smartphone, permitiendo asociar la presencia de una señal a una sala o zona determinada del museo. Aunque esta lógica es compatible con familias ampliamente difundidas como iBeacon o Eddystone~\cite{appleibeacon,googleeddystone,espressifibeacon}, el prototipo actual implementa un formato propio de \textit{service data} sobre BLE advertising, suficiente para codificar la identidad espacial requerida por la tesis sin incorporar una capa de complejidad adicional.

Para la solución propuesta, el dato relevante no es solamente el paquete emitido por el beacon, sino su interpretación como identificador lógico del espacio museográfico. Por ello, cada emisor debe estar descrito por un conjunto mínimo de campos: un \texttt{beacon\_id} interno, el protocolo utilizado, el payload identificador difundido, la sala o zona asociada y los parámetros básicos de emisión. En la implementación \texttt{iot-museiq}, este contrato se concreta mediante un \textit{service data} con UUID \texttt{0xA00A} que empaqueta \texttt{ROOM\_ID}, \texttt{BEACON\_NODE}, versión de firmware, potencia de transmisión y nivel de batería. A su vez, el identificador lógico legible se expone como \texttt{SALA\_x-Bnn}, lo que hace explícita la relación entre emisor físico y contexto museográfico. La Figura~\ref{fig:payload-ble-museiq} muestra este contrato desde la emisión del beacon hasta su lectura en la aplicación móvil.

\begin{figure}[!ht]
    \centering
    \includegraphics[width=0.96\textwidth]{payload_ble_museiq.png}
    \caption[Contrato BLE del beacon hacia la app]{Estructura del \textit{service data} BLE emitido por \texttt{iot-museiq} y su transformación en datos interpretables por la aplicación móvil. Fuente: elaboración propia.}
    \label{fig:payload-ble-museiq}
\end{figure}

\subsection{Inventario de beacons y mapeo espacial}

El reconocimiento de sala no puede depender únicamente de lecturas instantáneas de señal. Se requiere además un inventario previo que relacione cada beacon con su ubicación funcional. Dicho inventario debe contener, como mínimo, el identificador lógico del beacon, su identificador de broadcasting, la sala o zona asignada, una descripción breve de su ubicación física, la fecha de instalación y los parámetros relevantes de transmisión. Esta estructura hace posible mantener trazabilidad técnica, facilitar mantenimiento y evitar ambigüedades cuando se actualiza la infraestructura. En la maqueta operativa del proyecto, la convención \texttt{S1-M1}, \texttt{S1-M2} y \texttt{S1-M3} se usa como nombre corto de dispositivos de prueba, mientras que la aplicación móvil los normaliza a identificadores persistentes del tipo \texttt{SALA\_1-B01}, \texttt{SALA\_1-B02} y \texttt{SALA\_1-B03}. Esta correspondencia evidencia que el inventario no es un agregado documental, sino una dependencia efectiva para interpretar la señal observada.

En términos de diseño, este inventario es un dato de entrada tan importante como las lecturas capturadas por el teléfono. Sin él, la señal RSSI carecería de interpretación museográfica. Además, dado que el proyecto privilegia reconocimiento funcional antes que posicionamiento centimétrico, la relación \texttt{beacon $\rightarrow$ sala/zona} se convierte en el eje semántico del contexto espacial.

\subsection{Observaciones RSSI como señal de proximidad}

Durante la visita, el smartphone actúa como observador BLE y registra resultados de escaneo. Entre los campos más útiles para la solución propuesta se encuentran el identificador detectado, el valor RSSI y la marca temporal de observación. La documentación oficial de Android expone precisamente estos elementos en los resultados de escaneo BLE, lo que permite usar RSSI como aproximación a cercanía relativa sin requerir conexión persistente con el beacon~\cite{android2026scanresult}. En la aplicación móvil actual, estas observaciones se complementan con la dirección del dispositivo BLE, el \texttt{txPower} calibrado para estimación relativa y una marca de \texttt{lastSeen}, con el fin de distinguir beacons todavía activos de detecciones residuales.

No obstante, RSSI es una señal inherentemente ruidosa. La literatura sobre localización indoor advierte que la intensidad recibida puede verse afectada por orientación del teléfono, obstáculos, reflexiones y heterogeneidad del hardware receptor~\cite{leitch2023indoor,verde2023blemuseum}. Por esta razón, la solución propuesta no debe interpretar una sola lectura como evidencia suficiente de ubicación. La entrada funcional relevante corresponde, más bien, a una serie temporal corta de observaciones BLE que posteriormente podrá ser tratada mediante criterios de procesamiento definidos en el capítulo de diseño. El prototipo móvil ya refleja esta necesidad al aplicar ventanas cortas de RSSI y suavizado exponencial, lo que confirma que la unidad de entrada útil no es la lectura instantánea aislada, sino un patrón breve de observaciones sucesivas.

En consecuencia, cada observación BLE debe conservar, como mínimo, una marca temporal, el identificador del beacon detectado y el valor de intensidad de señal recibida. Cuando el entorno o la biblioteca empleada lo permitan, también conviene registrar la potencia de transmisión reportada. Este conjunto resulta suficiente para sostener el reconocimiento funcional por sala o zona, que es el objetivo realista del proyecto. La Figura~\ref{fig:mapeo-beacon-sala-museiq} sintetiza la relación entre emisión, inventario espacial y observación funcional.

\begin{figure}[!ht]
    \centering
    \includegraphics[width=0.94\textwidth]{mapeo_beacon_sala_museiq.png}
    \caption[Relación entre beacon, sala y observaciones BLE]{Relación entre infraestructura BLE, mapeo espacial y observaciones usadas como entrada funcional en la solución propuesta. Fuente: elaboración propia.}
    \label{fig:mapeo-beacon-sala-museiq}
\end{figure}

\section{Datos de entrada para orientación y apoyo contextual}

\subsection{Sensores del smartphone}

La capa de contexto de la solución propuesta no se limita a la proximidad BLE. El smartphone del visitante incorpora sensores que pueden complementar la interpretación del entorno, especialmente cuando el usuario permanece dentro de una misma sala y se requiere reducir ambigüedad respecto al foco probable de atención. Android documenta la posibilidad de derivar orientación combinando acelerómetro y sensor geomagnético, así como el uso de vectores de rotación cuando el dispositivo los soporta~\cite{android2026positionsensors}. En la aplicación \texttt{iot}, esta idea ya se plasma mediante lectura de acelerómetro, magnetómetro y podómetro, presentadas al usuario como indicadores de movimiento, rumbo aproximado y actividad durante el recorrido.

En este marco, la solución propuesta considera como entradas base las mediciones del acelerómetro y del magnetómetro, junto con sus respectivas marcas temporales. Cuando el hardware lo permita, también puede aprovechar sensores compuestos que entreguen orientación de forma más estable. El propósito no es construir navegación de alta precisión, sino disponer de una señal contextual complementaria que ayude a interpretar hacia dónde parece dirigirse el visitante. La implementación vigente calcula además un estado cualitativo de movimiento y un conteo de pasos nativo o estimado, de modo que la sesión de visita disponga no solo de orientación angular, sino también de una noción básica de desplazamiento.

\subsection{Variables derivadas de orientación}

A partir de las señales sensoriales base, el sistema puede derivar variables como azimuth, pitch y roll. Estas variables resultan útiles para discriminar orientación horizontal, inclinación y rotación lateral del dispositivo. En el contexto museístico, su valor reside en apoyar decisiones tales como priorizar el contenido de una vitrina o pieza ubicada frente al visitante, o reforzar la consistencia entre el historial reciente de proximidad y la dirección aparente de observación. Aunque el prototipo móvil hoy materializa esta capa sobre todo mediante \texttt{heading}, estado de movimiento y conteo de pasos, la formalización en términos de variables derivadas sigue siendo válida porque define el espacio de datos sobre el cual podrá evolucionar la inferencia contextual.

Aunque el cálculo detallado de estas variables corresponde al diseño de la solución, en este capítulo se fija su papel como entradas funcionales. También se reconoce que la calidad de estas mediciones depende del dispositivo y del entorno, por lo que deben acompañarse de metadatos de adquisición, tales como frecuencia de muestreo efectiva, disponibilidad de sensores y sincronización temporal con los eventos BLE.

\section{Datos de interacción del visitante}

\subsection{Consultas por voz y transcripción}

La interacción por voz es uno de los componentes diferenciales de la solución propuesta, tanto por accesibilidad como por naturalidad de uso. Desde el punto de vista de entradas, el sistema debe distinguir entre la señal acústica capturada por el micrófono y el dato funcional que realmente será procesado por la capa conversacional. En la práctica, la aplicación necesita registrar un identificador de consulta, una marca temporal, la transcripción producida por el servicio STT y, cuando la plataforma lo provea, un indicador de confianza o idioma detectado~\cite{android2026speech,wang2024voiceguide}. La app móvil ya dispone de captura de voz para dictado y transforma la interacción en una consulta textual que luego se envía al backend RAG.

La documentación de Android advierte que el reconocimiento de voz puede implicar el envío de audio a servicios remotos, hecho que introduce implicancias de privacidad y consumo de recursos~\cite{android2026speech}. En coherencia con ello, la solución propuesta debe privilegiar la minimización de datos: para la operación ordinaria, la transcripción y sus metadatos son más pertinentes que el almacenamiento persistente del audio crudo. Este criterio también refuerza la necesidad de definir reglas institucionales de tratamiento de datos personales en capítulos posteriores.

\subsection{Texto y eventos de interfaz}

Además de la voz, la solución propuesta requiere contemplar entradas equivalentes mediante texto escrito y eventos de interfaz. Un visitante puede formular consultas manualmente, seleccionar una pieza, activar una narración o elegir explícitamente una sala cuando el sistema necesite desambiguación. Estos eventos tienen valor contextual porque reducen incertidumbre y permiten al usuario corregir o complementar la interpretación automática del sistema.

Por ello, el esquema de entrada debe incorporar al menos \texttt{typed\_text}, \texttt{ui\_event}, \texttt{selected\_piece\_id} o campos equivalentes según la interacción ofrecida por la aplicación. Esta multimodalidad mantiene consistencia con la orientación inclusiva del proyecto y mejora robustez frente a ruido ambiental, diferencias de acento o limitaciones temporales del reconocimiento por voz. En el prototipo actual, el intercambio con \texttt{museRAG} ya se estructura mediante un \texttt{payload} que incluye \texttt{pregunta}, \texttt{museo}, \texttt{sala}, \texttt{obra}, \texttt{session\_id} y un bloque opcional \texttt{artwork\_context}. Este último aporta campos curatoriales inmediatos como \texttt{id}, \texttt{title}, \texttt{author}, \texttt{year}, \texttt{period}, \texttt{technique}, \texttt{summary}, \texttt{context}, \texttt{room\_relation}, \texttt{location\_hint} y \texttt{suggested\_questions}. En consecuencia, la entrada conversacional de la solución propuesta no es una frase aislada del usuario, sino una consulta contextualizada por el estado del recorrido, como se resume en la Figura~\ref{fig:payload-consulta-muserag}.

\begin{figure}[!ht]
    \centering
    \includegraphics[width=0.97\textwidth]{payload_consulta_muserag.png}
    \caption[Consulta contextual hacia MuseRAG]{Relación entre el \textit{payload} enviado por la app, las fuentes curatoriales indexadas y la respuesta contextual recuperada por \texttt{museRAG}. Fuente: elaboración propia.}
    \label{fig:payload-consulta-muserag}
\end{figure}

\subsection{Parámetros de síntesis de voz}

Aunque la síntesis de voz pertenece al dominio de salida, su operación depende de parámetros de entrada que deben estar previamente definidos. Entre ellos destacan idioma, voz seleccionada, tasa de locución, tono y criterios de segmentación de respuestas. Android proporciona soporte nativo para estas configuraciones mediante la API de \textit{Text-to-Speech}~\cite{android2026tts}. En consecuencia, la solución propuesta debe tratar estos parámetros como datos de configuración de sesión, pues afectan de manera directa la calidad de la experiencia interpretativa.

\section{Datos curatoriales para respuestas contextualizadas}

\subsection{Corpus institucional como fuente de conocimiento}

La capacidad conversacional de la solución propuesta no se sustenta únicamente en un modelo de lenguaje, sino en un corpus curatorial institucional que actúa como fuente principal de conocimiento. Este corpus debe incluir, como mínimo, fichas de piezas, textos de sala, descripciones interpretativas, materiales educativos y contenidos complementarios relevantes para la visita. La función del corpus es proveer evidencia recuperable para que la respuesta generada se mantenga alineada con la narrativa museográfica y reduzca la probabilidad de alucinación~\cite{lewis2020rag,karpukhin2020dpr,bechard2024hallucination}. En la base actual del proyecto, dicho corpus ya se construye a partir de tres fuentes explícitas: un PDF institucional sobre Tumbas Reales de Sipán, el archivo \texttt{museum.json} con narrativas por sala y zona, y el archivo \texttt{datos.ts} que define obras, textos y relaciones del recorrido.

En ese sentido, el corpus no debe concebirse como una colección informal de textos, sino como un activo institucional gobernado. Su calidad influye directamente en la precisión semántica del sistema y en la capacidad de justificar las respuestas ofrecidas al visitante. Esta necesidad de gobernanza resulta además coherente con la orientación de UNESCO sobre el uso responsable de inteligencia artificial en contextos culturales, donde la tecnología debe reforzar y no diluir la autoridad del conocimiento patrimonial~\cite{unescoaiculture2025}.

\subsection{Metadatos mínimos para recuperación}

Para que el corpus sea utilizable por una estrategia de recuperación aumentada, cada unidad documental debe incorporar metadatos consistentes. Como mínimo, la solución propuesta requiere identificar la pieza o sala asociada, el idioma del contenido, el tipo documental, la fuente institucional, la versión del contenido y los derechos o restricciones de uso cuando existan recursos multimedia asociados. Estos campos permiten filtrar, recuperar y trazar evidencia de manera más confiable durante la conversación. En \texttt{museRAG}, esta lógica ya se expresa mediante metadatos como \texttt{museum\_id}, \texttt{room\_id}, \texttt{artwork\_id}, \texttt{kind}, \texttt{source} e incluso \texttt{image\_url} para vincular texto recuperado con apoyo visual en la interfaz móvil.

La investigación revisada para este capítulo sugiere además que la organización del conocimiento patrimonial se beneficia de referencias a estándares de interoperabilidad y catalogación, como CIDOC CRM, LIDO y Dublin Core, aun cuando el proyecto no implemente exhaustivamente dichos modelos~\cite{cidoccrm,lido,dublincore}. Para la solución propuesta, la lección práctica no es adoptar una ontología compleja desde el prototipo, sino asumir el principio de metadatos mínimos consistentes y estables.

\subsection{Segmentación y unidad recuperable}

Otro aspecto clave es la definición de la unidad recuperable que alimentará la capa RAG. Un documento curatorial completo puede ser demasiado amplio o heterogéneo para su recuperación directa; por ello, el sistema debe prever una segmentación en fragmentos o \textit{chunks} que conserven relación explícita con su documento de origen, la sala, la pieza y el idioma correspondiente. Esta estructura favorece recuperación más precisa y hace posible citar o rastrear la evidencia que sustenta una respuesta. El backend implementado ya sigue esta decisión al indexar fragmentos provenientes del PDF, las narrativas de sala y las fichas curatoriales, y al devolverlos como \texttt{SourceSnippet} con puntaje, metadatos y, cuando corresponde, imagen asociada.

Así, la solución propuesta requiere al menos las entidades lógicas \texttt{documento curatorial} y \texttt{chunk recuperable}, relacionadas entre sí mediante identificadores estables. Esta decisión no pertenece todavía a la implementación detallada del índice, pero sí constituye una definición necesaria de los datos que alimentarán la recuperación posterior.

\section{Organización lógica e integridad de los datos}

\subsection{Identificadores y relaciones}

La integración de señales físicas, eventos de interacción y conocimiento curatorial exige un esquema coherente de identificadores internos. La solución propuesta debe manejar, como mínimo, identificadores estables para museo, sala, zona, beacon, pieza, documento curatorial, chunk recuperable y sesión de interacción. Sin este principio, no sería posible unir observaciones BLE con contexto museográfico, ni filtrar el corpus de acuerdo con la ubicación o la consulta del visitante.

Sobre esta base, puede plantearse un modelo lógico mínimo compuesto por entidades como \texttt{Sala}, \texttt{Beacon}, \texttt{Pieza}, \texttt{DocumentoCuratorial}, \texttt{Chunk}, \texttt{EventoBLE}, \texttt{EventoSensor} y \texttt{Consulta}. Lo importante en esta etapa no es fijar un esquema físico de base de datos ni una arquitectura de persistencia, sino establecer las relaciones imprescindibles entre contexto espacial, interacción y conocimiento institucional. Sin embargo, el prototipo ya ofrece una traducción concreta de este modelo: la app local mantiene tablas equivalentes a \texttt{museum\_profile}, \texttt{rooms}, \texttt{artworks}, \texttt{route\_steps}, \texttt{faq}, \texttt{voice\_prompts} y \texttt{permissions\_catalog}; el backend RAG maneja entidades de consulta, respuesta y fuentes recuperadas; y la capa IoT materializa \texttt{Beacon} y observaciones BLE mediante identificadores de sala y nodo. La Figura~\ref{fig:modelo-logico-datos-museiq} condensa esta estructura mínima y sus relaciones principales.

\begin{figure}[!ht]
    \centering
    \includegraphics[width=0.96\textwidth]{modelo_logico_datos_museiq.png}
    \caption[Modelo lógico mínimo de datos de entrada]{Modelo lógico mínimo de entidades y relaciones para los datos de entrada de la solución propuesta. Fuente: elaboración propia.}
    \label{fig:modelo-logico-datos-museiq}
\end{figure}

\subsection{Criterios mínimos de calidad}

La tesis también debe explicitar que la utilidad de los datos de entrada depende de requisitos mínimos de calidad. En esta propuesta, dichos requisitos son cuatro. El primero es consistencia: el inventario de beacons no puede contener asignaciones contradictorias entre emisor y sala. El segundo es completitud: cada documento curatorial debe contar al menos con identificador, idioma, fuente y relación con una sala o pieza. El tercero es actualidad: el corpus debe mantenerse versionado para que la recuperación privilegie contenido vigente. El cuarto es trazabilidad: cada respuesta contextual debe poder vincularse con evidencia institucional recuperada.

Estos criterios se alinean con marcos generales de calidad de datos, como ISO/IEC 25012, que enfatizan dimensiones verificables para datos estructurados dentro de sistemas computacionales~\cite{iso25012}. En la solución propuesta, su formulación es especialmente importante porque la percepción de calidad por parte del visitante depende tanto de la infraestructura técnica como de la confiabilidad del contenido ofrecido. En términos de implementación, ello se refleja en necesidades muy concretas: que el \texttt{ROOM\_ID} emitido por el ESP32 coincida con la sala registrada en la app, que las obras activas del recorrido mantengan identificadores estables entre \texttt{datos.ts}, almacenamiento local y backend RAG, y que toda respuesta conversacional conserve referencias a los fragmentos fuente utilizados para generarla.

\subsection{Seguridad y pertinencia de las entradas}

Finalmente, no todos los datos de entrada deben tratarse con el mismo grado de confianza. El corpus curatorial debe provenir de fuentes institucionales controladas, mientras que las consultas de usuario en voz o texto deben asumirse como entradas no confiables desde la perspectiva del sistema conversacional. Esta distinción es relevante porque aplicaciones basadas en modelos de lenguaje y RAG pueden verse afectadas por intentos de desviar instrucciones o introducir contenido inapropiado a través de las propias entradas del usuario~\cite{owaspllm}.

Del mismo modo, las transcripciones, preferencias y registros de sesión pueden constituir datos personales o potencialmente sensibles, por lo que su tratamiento debe alinearse con criterios de minimización y con el marco normativo peruano de protección de datos personales~\cite{ley29733}. Aunque el desarrollo normativo se abordará con mayor detalle en capítulos posteriores, en este punto ya puede establecerse que la pertinencia del sistema depende también de cómo gobierna y limita sus datos de entrada.

\section{Síntesis del capítulo}

El diseño de datos de entrada de la solución propuesta se apoya en una combinación de señales de proximidad BLE, variables de orientación del smartphone, consultas multimodales del visitante y un corpus curatorial institucional estructurado. Cada uno de estos insumos cumple una función específica dentro del sistema: reconocer sala o zona probable, reducir ambigüedad contextual, capturar intención del visitante y sustentar respuestas conversacionales con evidencia documental. El acceso a los repositorios base del proyecto confirma que esta estructura ya no es solamente hipotética: existe un prototipo BLE con contrato explícito de advertising, una aplicación móvil con almacenamiento y sensado local, y un backend RAG que recibe contexto de museo, sala y obra para recuperar fuentes y responder. En conjunto, las Figuras~\ref{fig:flujo-datos-museiq-ilustrativo}, \ref{fig:payload-ble-museiq}, \ref{fig:mapeo-beacon-sala-museiq}, \ref{fig:payload-consulta-muserag} y \ref{fig:modelo-logico-datos-museiq} permiten visualizar esta continuidad desde la captura de señales hasta la estructuración lógica de los insumos.

En consecuencia, este capítulo establece el contrato de datos que hará posible el diseño del sistema en el capítulo siguiente. La arquitectura y el despliegue de la solución deberán partir de estas definiciones para integrar adquisición sensorial, manejo de interacción y recuperación de conocimiento de manera coherente, trazable y compatible con el entorno museístico en el que se propone operar el sistema.

